\documentclass[a4paper,12pt,oneside]{amsbook}
\usepackage[T1]{fontenc}
\usepackage{textcomp,hyperref}
%\usepackage{euler}
%\usepackage{garamond}

\usepackage{newcent}
%\fontfamily{fxm}\selectfont
\usepackage{fancyhdr}
\pagestyle{fancy}
% Detta paket skoter grafiken
\usepackage[dvips]{graphicx}
%\usepackage[pdftex]{graphicx}
\usepackage{tabularx}
\usepackage{amsthm}
\usepackage{amsmath}
\usepackage[swedish]{babel}
% For bortkommentering
\usepackage{verbatim}
% For att slippa indentering av Sections och SubSections
\usepackage{indentfirst}
\usepackage{hyperref}

% Har stalls paragrafindentering och avstand mellan stycken och rader mm
\setlength{\parindent}{0pt}
\setlength{\parskip}{6pt}
\setlength{\baselineskip}{12pt}
% Har definieras sidan
\setlength{\voffset}{-20mm}
\setlength{\hoffset}{-5.5mm}

\setlength{\topmargin}{0mm}
\setlength{\headheight}{10mm}
\setlength{\headsep}{8mm}
\setlength{\footskip}{8mm}

\setlength{\evensidemargin}{0mm}
\setlength{\oddsidemargin}{0mm}

\setlength{\marginparsep}{0mm}
\setlength{\marginparwidth}{0mm}

\setlength{\textwidth}{170mm}
\setlength{\textheight}{240mm}
\setlength{\paperwidth}{210mm}
\setlength{\paperheight}{297mm}
\setlength{\headwidth}{170mm}
% Huvud och Fot
%\fancypagestyle{}{
\fancyhf{}\fancyhead[c]{\small{\textit{Svar och rättningsanvisningar
      till Programmeringsolympiadens skolkval 2024}}}
\fancyfoot[c]{\thepage}
\renewcommand{\headrulewidth}{0.4pt}
\renewcommand{\footrulewidth}{0.4pt}
\newenvironment{lista}
{\begin{itemize}
\setlength{\parindent}{0pt}
\setlength{\itemsep}{6pt}}
{\end{itemize}}
% Uppgifter
\newcounter{probnr}
\newenvironment{tal}{%
\begin{list}
%{\textbf{\arabic{section}.\arabic{probnr}}} {\usecounter{probnr}
{\textbf{\arabic{probnr}}} {\usecounter{probnr}
\setlength{\leftmargin}{0mm}
\setlength{\rightmargin}{0mm}
\setlength{\labelwidth}{-1mm}
\setlength{\labelsep}{1mm}}
\setlength{\itemsep}{6pt}
}{\end{list}}
% Definition av avsnitt: FAKTA, EXEMPEL, PASTAENDE
\newtheorem{fakta}{Fakta}
\newtheorem{exempel}{Exempel}
\newtheorem{problem}{Problem}
%
\newtheoremstyle{test}% NAME
{20pt}      % ABOVESPACE
{10pt}      % BELOWSPACE
{\sffamily} % BODYFONT
{0pt}       % INDENT
{\scshape}  % HEADFONT
{}          % HEADPUNCT
{\newline}  % HEADSPACE
{}          % CUSTOM-HEAD-SPEC

\theoremstyle{test}
\newtheorem{program}{Program}
\newcommand{\sv}[1]{\textsc{#1}}            % Sma versaler
\newcommand{\fe}[1]{\textbf{#1}}            % FET
\newcommand{\ku}[1]{\textit{#1}}            % KURSIV
\newcommand{\cu}[1]{\texttt{#1}}            % Courier
\newcommand{\sk}[1]{\texttt{#1}}            % Courier
\newcommand{\rubrik}[1]{\begin{center}\sf\huge{#1}\normalsize\rm\end{center}}
\begin{document}
%\DeclareGraphicsExtensions{.jpg,.pdf,.mps,.png,.eps}


\specialsection*{Svar och rättningsanvisningar}
\thispagestyle{fancy}
\lhead{}
\begin{itemize}
%\setlength\itemsep{0.2cm}
\item Läs igenom \textit{tävlingsreglerna}.
\item Programmen tas i tur och ordning in i editorn och kompileras.
Uppstår kompileringsfel betraktas programmet som felaktigt och lösningen
ges $0$ poäng.
\item Programmet körs med givna indata enligt nedan. Alternativt, om eleven förberett programmet för det, kan de bifogade indatafilerna användas istället.
\item Varje testfall
  med korrekt svar ger 1 poäng.
\item
Totalt kan man
  alltså få 5 poäng för varje uppgift.
\item Om exekveringstiden för ett testexempel, kört på en modern dator,
\textit{överskrider 3 sekunder} betraktas körningen av testexemplet som felaktigt.
\item Det kan vara viktigt att programmet körs i en miljö liknande den som
programmet utvecklats i, samma version av kompilator eller
interpretator.
\item Vid problem i samband med rättningen är det viktigt att det sunda
förnuftet får råda!
\item Ett förslag till rättningsprocedur kan vara att låta eleven
sitta vid datorn.
\end{itemize}

\vspace{1cm}

\subsection*{Uppgift 1 -- Reduplikation}
~\\
{\tt 
\begin{tabular}{||l||c|c||c||}\hline\hline
& \multicolumn{2}{c||}{\fe{Indata}} & \multicolumn{1}{c||}{\fe{Utdata}} \\ 
& Strängen & Repetitioner &  \\ \hline \hline
\fe{Test 1} & hej &1&hej\\ \hline
\fe{Test 2} & lol &4&lollollollol \\ \hline
\fe{Test 3} & naj &5&najnajnajnajnaj \\ \hline
\fe{Test 4} & python&3&pythonpythonpython \\ \hline
\fe{Test 5} & ha &8&hahahahahahahaha\\ \hline\hline
\end{tabular}
}
%\vspace{1cm}


\subsection*{Uppgift 2 -- Tre i rad}
~\\
{\tt 
\begin{tabular}{||l||l||c||}\hline\hline
& {\fe{Indata}} & \fe{Utdata} \\ \hline \hline
\fe{Test 1} & 13 & 1\\ \hline
\fe{Test 2} & 69 & 3 \\ \hline
\fe{Test 3} & 992 & 9 \\ \hline
\fe{Test 4} & 91081 & 44\\ \hline
\fe{Test 5} & 907038264 & 966\\ \hline\hline
\end{tabular}
}

\subsection*{Uppgift 3 -- Meddelande}
~\\
{\tt 
\begin{tabular}{||l||c|c|l||l||}\hline\hline
& \multicolumn{3}{c||}{\fe{Indata}} & \fe{Utdata} \\
& $n$ & $m$ & Rader & \\ \hline \hline
\fe{Test 1} & 1 & 6 &	.he..j &	hej	\\ \hline
\fe{Test 2} & 6 & 5 & r....	&	ratvinklig	\\
&&&	a....	&	\\
&&&	t....	&	\\ 
&&&	v....	&	\\ 
&&&	i....	&	\\ 
&&&	nklig	&	\\ \hline
\fe{Test 3} & 6 & 6 & tab...	&	tabortpunkt	\\
&&&	..or..	&	\\
&&&	...t..	&	\\
&&&	...pu.	&	\\
&&&	....nk	&	\\
&&&	.....t	&	\\ \hline
\fe{Test 4} & 5 & 6 &	......	&	center	\\
&&&	......	&	\\ 
&&&	center	&	\\
&&&	......	&	\\ 
&&&	......	&	\\ \hline
\fe{Test 5} & 6 & 6 &	.s....	&	svar	\\
&&&	......	& \\
&&&	...v..	& \\
&&&	.....a	& \\
&&&	......	& \\
&&&	.....r	& \\\hline \hline
\end{tabular}
}

\subsection*{Uppgift 4 -- Klippa protein}
~\\
{\tt 
\begin{tabular}{||l||c|c||c||}\hline\hline
& \multicolumn{2}{c||}{\fe{Indata}} & \multicolumn{1}{c||}{\fe{Utdata}} \\ 
& Strängen & $k$ &  \\ \hline \hline
\fe{Test 1} & thequickbrownfoxjumpsoverthe &1&25\\ \hline
\fe{Test 2} & abbbababbbbabababbabbbabbaaababb &5&5 \\ \hline
\fe{Test 3} & skolkvalprogrammeringsolympiad &2&14 \\ \hline
\fe{Test 4} & dennameningartrettiotvabokstaver &2&14 \\ \hline
\fe{Test 5} & hejsandettaarettstarkttestfall &7&4\\ \hline\hline
\end{tabular}
}
\vspace{1cm}

\subsection*{Uppgift 5 -- Miniräknaren}
~\\
{\tt 
\begin{tabular}{||l||c|c||c||}\hline\hline
& \multicolumn{2}{c||}{\fe{Indata}} & \multicolumn{1}{c||}{\fe{Utdata}} \\ 
& n & m &  \\ \hline \hline
\fe{Test 1} & 9 & 3 & 3 \\ \hline
\fe{Test 2} & 54581 & 2 & 24 \\ \hline
\fe{Test 3} & 99896 & 7 & 11 \\ \hline
\fe{Test 4} & 384463392 & 7 & 20 \\ \hline
\fe{Test 5} & 488339432 & 9 & 16\\ \hline\hline
\end{tabular}
}
\vspace{1cm}

\subsection*{Uppgift 6 -- Pianot}
~\\
{\tt 
\begin{tabular}{||l||l||c||}\hline\hline
& {\fe{Indata}} & \fe{Utdata} \\ \hline \hline
\fe{Test 1} & dettafallenkelt & 17 \\ \hline
\fe{Test 2} & hoppaslosningenklarardettafall & 58 \\ \hline
\fe{Test 3} & thisistestcasenumberthree & 41 \\ \hline
\fe{Test 4} & clownassociationenslemurfarm & 61\\ \hline
\fe{Test 5} & yxskaftmakareombudsmannen & 48\\ \hline\hline
\end{tabular}
}




\end{document}




\end{document}


